\section{Discussion}
\label{sec:discussion}

\subsection{Interpretation of the main claim}

The experiments support a focused claim: in the evaluated simulator, secrecy-aware synchronization should be optimized jointly with UAV motion, source power, and cooperative jamming. The evidence is strongest in the calibrated stress scenario, where stale or low-quality twin information has a visible effect on outage. In that setting, \texttt{rollout\_joint} improves both average secrecy rate and outage relative to periodic, risk-triggered, and margin-triggered rules.

The base scenario should be interpreted more cautiously. Although \texttt{rollout\_joint} has the highest average secrecy rate, its outage probability is the same as \texttt{periodic}. This means the base case demonstrates competitiveness and modest secrecy gain rather than a major reliability improvement. The hard scenario sits between these two regimes. This pattern is scientifically useful because it shows that the proposed method is most valuable when synchronization quality and outage risk are actually limiting factors.

\subsection{Why the oracle and ablation results matter}

The oracle-state rollout establishes that better eavesdropper state information can substantially reduce outage. The gap between \texttt{oracle\_sync} and \texttt{rollout\_joint} therefore quantifies remaining room for improvement in twin accuracy, synchronization reliability, and control approximation. The \texttt{rollout\_fixed\_periodic} result is also important: it uses rollout search but removes adaptive synchronization. Its weaker performance indicates that the proposed gain is not only caused by better motion and power search. The \texttt{rollout\_no\_sync} and \texttt{no\_twin} results further show that neither stale-state rollout nor twin-free control is sufficient in the stress scenario.

\subsection{Scope of the certificate}

The certificate improves the scientific structure of the paper because it prevents the controller from relying blindly on twin-predicted secrecy. The holdout validation results show that the fitted upper bound is conservative on the evaluated trajectories. Nevertheless, the certificate is empirical and policy-dependent. It should be described as an in-policy risk estimator, not as a formal proof of safety for all possible eavesdropper behaviors. A stronger future version would test distribution shift explicitly, for example by changing Eve mobility, channel parameters, synchronization failure probability, and user layouts after calibration.

\subsection{Limitations}

The study has four main limitations. First, the simulator abstracts many real-world effects, including sensing errors, UAV dynamics, regulatory constraints, and hardware-level communication overhead. Second, the rollout controller is computationally expensive, requiring about 1.2--1.3 seconds per slot in the main experiments. Third, the PPO baseline uses a limited 200-episode training budget and should be treated only as a lightweight comparison. Fourth, the exact-DP experiment is a sanity check on a reduced model, not a global optimality certificate for the paper-scale scenarios.

These limitations do not invalidate the main result, but they define the appropriate strength of the paper's claims. The work is best positioned as a simulation-based communication-system study showing the value of jointly optimizing digital-twin synchronization and physical-layer security control.

\subsection{Future work}

Future work should reduce runtime through parallel rollout evaluation, learned value approximators, candidate pruning, and adaptive horizon selection. The certificate should be extended to explicit distribution-shift validation and possibly to stronger conformal-style coverage conditions. Larger networks with more UAVs, multiple eavesdroppers, and richer channel dynamics should also be studied. Finally, reinforcement-learning baselines should be revisited with larger training budgets and architectures designed specifically for the coupled synchronization-control action space.
